\documentclass{report}
\usepackage{amsfonts, amsmath}
\newcommand\R{{\mathbb R}}
\newcommand\Q{{\mathbb Q}}
\newcommand\N{{\mathbb N}}
\newcommand{\vf}{\varphi}
\pagestyle{empty}
\begin{document}
\hfuzz 12pt
\large
\centerline{\Huge \bf Analisi Matematica II}
\vskip.2cm
\centerline{\Huge Correzione Secondo Appello}

\vskip1cm
\begin{enumerate}
\item $g(x+2\pi)=f(3x+6\pi)=f(3x)=g(x)$, dunque $g$ \`e periodica.
\[
\begin{split}
a_n(f)=&\frac{1}{2\pi}\int_0^{2\pi}f(x)\cos
nx\;dx
=\frac{1}{6\pi}\int_0^{6\pi}g(x/3)\cos nx\;dx \\
&=\frac{1}{2\pi}\int_0^{2\pi}g(x)\cos 3nx\;dx=a_{3n}(g).
\end{split}
\]
e, analogamente, $b_n(f)=b_{3n}(g)$.

Se $f=g$ allora segue $a_n(g)=a_n(f)=a_{3n}(g)$. Questo non dice
nulla per $n=0$ ma per $n\neq 0$ implica $a_n(g)=a_{3^kn}(g)$ per ogni
$k>0$. Ma poich\`e la funzione $g$ \`e derivabile segue che
$\lim_{n\to\infty}a_n(g)=\lim_{n\to\infty} b_n(g)=0$. Dunque
$a_n(g)=b_n(0)=0$ per ongi $n\neq0$. Finalmente, poich\`e la funzione
\`e derivabile due volte, ne segue che \`e uguale alla sua serie di
Fourier, ma ci\`o significa $g(x)=a_0(g)$, cio\`e $g$ , e dunque anche
$f$, \`e costante.

\item Poich\`e la funxione \`e derivabile segue da Lagrange che
$f(x)=f(0)+f'(\xi)x$
per qualche punto $\xi\in (0,x)$. Dunque
\[
n\int^{\frac 1n}_0f(x)=f(0)+n\int^{\frac 1n}_0f'(\xi)x.
\]
Dunque 
\[
\left|n\int^{\frac 1n}_0f(x)-f(0)\right|\leq \frac L{2n}
\]
da dui segue
\[
\lim_{n\to\infty}n\int^{\frac 1n}_0f(x)=f(0).
\]
Lo stesso risultato poteva essere ottenuto integrando per parti
($int^{\frac 1n}_0f(x)dx=f(1/n)-int^{\frac 1n}_0f'(x)xdx$) e poi
argomentando come sopra, oppure (per gli amanti del genere) usando
L'Hopital.

\item Poich\`e
\[
2e^x-1-(x+1)^2=2[e^x-1-x-\frac 12x^2]=2[\frac 16e^\xi x^3]
\]
dove abbiamo usato la furmula di Taylor al terzo ordine per
l'esponenziale, segue che il limite \`e uguale a $1/3$.
\item La soluzione della equazione differenziale \`e
$x(t)=1-e^{-\gamma t}$, dunque il limite cercato \`e uguale ad uno.
\item Chiaramente
\[
A(\lambda)=\int_0^\lambda p_\lambda(x)dx=\lambda^{-1}\int_0^1x(1-x)dx
+\lambda=\frac 1{6\lambda}+\lambda.
\]
Poch\`e $A(\lambda)$ tende all'infinito sia per $\lambda$ tendente a
zero che all'infinito, segue che deve esiste un minimo e per trovarlo
basta vedere dove la derivata di $A$ rispetto a $\lambda$ si annulla.

\item Il settore circolare si puo vedere come unione di un triangolo
di base $r\cos\alpha$ e altezza $r\sin \alpha$ pi\`u il rimanente, cio\`e
l'area cercata \`e data da 
\[
\frac
{r^2}2\cos\alpha\sin\alpha+\int^{1}_{r\cos\alpha}\sqrt{r^2-x^2}dx
\]
Usando il cambio di variabile nell'integrale $z= r\cos x$ si ottine il
risultato $\frac{r^2\alpha}2$.\footnote{\`E interessante notare che la
formula ha una stretta analogia con quella dell'area del trinagolo:
lunghezza della base per ``altezza'' diviso due, comunque una tale
analogia non costituisce una dimostrazione. Allo stesso modo non \`e
immediatamente una dimostrazione ottenere il risultato dicendo che
l'area cercata \`e $2\pi/\alpha$  parte dell'area del cerchio. Questo
\`e ovvio se $2\pi/\alpha$ \`e un numero razionale (la dimostrazione
sono i soli spicchi di torta che ci perseguitano dalle elementari) ma
se \`e irrazionale per avere una dimostrazione occorre un argomento
non banale (bisogna approssimare i reali con i razionali, come al
solito).} 
\end{enumerate}
\end{document}
